% GENERATED FILE. Edit canonical lessons, then run npm run generate:books.
% canonical-source-hash: fnv1a64:a95004e2bc399296
% canonical-lessons: GU-C137-garib, GU-C137-juvan, GU-C137-majbut, GU-C137-nabalu, GU-C137-jaruri

\chapter{Poor, Young, Strong, Weak, Necessary}
\label{ch:gu-poor-young-strong-weak-necessary}
\hlchaptermodality{137}

\begin{chapteropening}
\textbf{By the end of this chapter:} \emph{I can say poor, young, strong, weak and necessary.}

Complete the last lesson of chapter 137: i can say poor, young, strong, weak and necessary.
\end{chapteropening}

\section[garīb]{\gu{ગરીબ} (garīb) — poor}

\label{lesson:GU-C137-garib}



\begin{quote}
Before the new one: say the Gujarati for false, then the Gujarati for rich.
\end{quote}

\subsection*{You'll want to know: \gu{ગરીબ}}
\textbf{\gu{ગરીબ}} (\emph{garīb}) — \textquotedblleft{}poor\textquotedblright{}.

\begin{grammarlens}[title={What you've built}]
One more word for this chapter.
\end{grammarlens}

\subsection*{Your turn}
\begin{itemize}
\raggedright
  \item \emph{Say it:} \emph{garīb}
  \item \emph{Say it:} \emph{garīb}, once more
  \item \emph{Say it:} say \emph{garīb}
  \item \emph{Recall:} say the Gujarati for false, then the Gujarati for rich, then say \emph{garīb} again
\end{itemize}

\begin{tcolorbox}[breakable,colback=teal!4,colframe=teal!35!black,title={Before you move on}]
What does \gu{ગરીબ} mean? (\textquotedblleft{}Poor\textquotedblright{}.) Say it once more.
\end{tcolorbox}
\section[juvān]{\gu{જુવાન} (juvān) — young}

\label{lesson:GU-C137-juvan}



\begin{quote}
Before the new one: say the Gujarati for rich, then the Gujarati for poor.
\end{quote}

\subsection*{You'll want to know: \gu{જુવાન}}
\textbf{\gu{જુવાન}} (\emph{juvān}) — \textquotedblleft{}young\textquotedblright{}.

\begin{grammarlens}[title={What you've built}]
One more word for this chapter.
\end{grammarlens}

\subsection*{Your turn}
\begin{itemize}
\raggedright
  \item \emph{Say it:} \emph{juvān}
  \item \emph{Say it:} \emph{juvān}, once more
  \item \emph{Say it:} say \emph{juvān}
  \item \emph{Recall:} say the Gujarati for rich, then the Gujarati for poor, then say \emph{juvān} again
\end{itemize}

\begin{tcolorbox}[breakable,colback=teal!4,colframe=teal!35!black,title={Before you move on}]
What does \gu{જુવાન} mean? (\textquotedblleft{}Young\textquotedblright{}.) Say it once more.
\end{tcolorbox}
\section[majbūt]{\gu{મજબૂત} (majbūt) — strong}

\label{lesson:GU-C137-majbut}



\begin{quote}
Before the new one: say the Gujarati for poor, then the Gujarati for young.
\end{quote}

\subsection*{You'll want to know: \gu{મજબૂત}}
\textbf{\gu{મજબૂત}} (\emph{majbūt}) — \textquotedblleft{}strong\textquotedblright{}.

\begin{grammarlens}[title={What you've built}]
One more word for this chapter.
\end{grammarlens}

\subsection*{Your turn}
\begin{itemize}
\raggedright
  \item \emph{Say it:} \emph{majbūt}
  \item \emph{Say it:} \emph{majbūt}, once more
  \item \emph{Say it:} say \emph{majbūt}
  \item \emph{Recall:} say the Gujarati for poor, then the Gujarati for young, then say \emph{majbūt} again
\end{itemize}

\begin{tcolorbox}[breakable,colback=teal!4,colframe=teal!35!black,title={Before you move on}]
What does \gu{મજબૂત} mean? (\textquotedblleft{}Strong\textquotedblright{}.) Say it once more.
\end{tcolorbox}
\section[nabaḷũ]{\gu{નબળું} (nabaḷũ) — weak}

\label{lesson:GU-C137-nabalu}



\begin{quote}
Before the new one: say the Gujarati for young, then the Gujarati for strong.
\end{quote}

\subsection*{You'll want to know: \gu{નબળું}}
\textbf{\gu{નબળું}} (\emph{nabaḷũ}) — \textquotedblleft{}weak\textquotedblright{}.

\begin{grammarlens}[title={What you've built}]
One more word for this chapter.
\end{grammarlens}

\subsection*{Your turn}
\begin{itemize}
\raggedright
  \item \emph{Say it:} \emph{nabaḷũ}
  \item \emph{Say it:} \emph{nabaḷũ}, once more
  \item \emph{Say it:} say \emph{nabaḷũ}
  \item \emph{Recall:} say the Gujarati for young, then the Gujarati for strong, then say \emph{nabaḷũ} again
\end{itemize}

\begin{tcolorbox}[breakable,colback=teal!4,colframe=teal!35!black,title={Before you move on}]
What does \gu{નબળું} mean? (\textquotedblleft{}Weak\textquotedblright{}.) Say it once more.
\end{tcolorbox}
\section[jarūrī]{\gu{જરૂરી} (jarūrī) — necessary}

\label{lesson:GU-C137-jaruri}



\begin{quote}
Before the new one: say the Gujarati for strong, then the Gujarati for weak.
\end{quote}

\subsection*{You'll want to know: \gu{જરૂરી}}
\textbf{\gu{જરૂરી}} (\emph{jarūrī}) — \textquotedblleft{}necessary\textquotedblright{}.

\begin{grammarlens}[title={What you've built}]
One more word for this chapter.
\end{grammarlens}

\subsection*{Your turn}
\begin{itemize}
\raggedright
  \item \emph{Say it:} \emph{jarūrī}
  \item \emph{Say it:} \emph{jarūrī}, once more
  \item \emph{Say it:} say \emph{jarūrī}
  \item \emph{Recall:} say the Gujarati for strong, then the Gujarati for weak, then say \emph{jarūrī} again
\end{itemize}

\begin{tcolorbox}[breakable,colback=teal!4,colframe=teal!35!black,title={Before you move on}]
What does \gu{જરૂરી} mean? (\textquotedblleft{}Necessary\textquotedblright{}.) Say it once more.
\end{tcolorbox}
